\documentclass[11pt]{article}

\usepackage[margin=1.25in]{geometry}
\usepackage[T1]{fontenc}
\usepackage{amsmath,amssymb,amsthm,mathtools}
\usepackage{bm}
\usepackage{enumitem}
\usepackage{microtype}
\usepackage{xcolor}
\usepackage[
  colorlinks=true,
  linkcolor=blue!55!black,
  citecolor=blue!55!black,
  urlcolor=blue!55!black
]{hyperref}
\usepackage[noabbrev,nameinlink]{cleveref}

\newtheorem{theorem}{Theorem}[section]
\newtheorem{lemma}[theorem]{Lemma}
\newtheorem{proposition}[theorem]{Proposition}
\newtheorem{corollary}[theorem]{Corollary}
\newtheorem{observation}[theorem]{Observation}
\newtheorem{definition}[theorem]{Definition}
\theoremstyle{remark}
\newtheorem{remark}[theorem]{Remark}

\newcommand{\R}{\mathbb{R}}
\newcommand{\one}{\bm{1}}
\newcommand{\nnz}{\operatorname{nnz}}
\newcommand{\OPT}{\operatorname{OPT}}

\title{A reduction from linear programming to positive linear programming}
\author{}
\date{July 19, 2026}

\begin{document}

\maketitle

\begin{abstract}
We give a deterministic polynomial-time reduction from the
threshold-decision problem for binary-encoded rational linear programs
to an exact positive packing-LP threshold problem.  Given
$A\in\mathbb{Q}^{m\times n}$, $b\in\mathbb{Q}^m$,
$c\in\mathbb{Q}^n$, and $\tau\in\mathbb{Q}$, the source question is
whether there is an $x\in\R^n$ such that
\[
  Ax\leq b,
  \qquad
  c^\top x\geq\tau.
\]
The reduction maps this question to whether the optimum of a positive
packing linear program reaches a specified threshold.  It first
computes a box of polynomial encoding length that contains a feasible
point whenever one exists.  The output has at most $2n$ variables and
at most $m+n+1$ constraints, and its constraint matrix $P$ satisfies
$\nnz(P)\leq\nnz(A)+\nnz(c)+2n$.  Every output coefficient has
polynomial encoding length.  Once the rational rows have been
normalized and the box computed, the sparse core uses
$O(\nnz(A)+\nnz(c)+m+n)$ exact arithmetic operations and output writes.
An explicit packing point at the threshold decodes to a source feasible
point in $O(n)$ time.

The reduction writes the objective threshold as one additional linear
inequality, splits signed variables into nonnegative pairs, shifts each
signed row using the computed box bound, and uses the positive objective
to force the bounding inequalities to be tight.
One motivation is the quantitative reduction of Kyng, Wang, and Zhang
linking accurate packing-LP solving to linear-system solving.  The
present construction is deliberately exact rather than approximation-preserving:
infeasible source instances can map to packing optima arbitrarily close
to the threshold.  It therefore does not by itself transfer running
times for multiplicative positive-LP approximation.
\end{abstract}

\begin{quote}
\small\noindent\textbf{Note.} This paper was fully AI generated with
ChatGPT 5.5 Pro from prompts by Xinyuan Cao, Hanxi Lin, and Richard Peng.
The present repository version, including this disclosure note, was
audited and revised by GPT-5 in the Codex harness.
\end{quote}

\section{Introduction}
\label{sec:introduction}

General linear programs may have signed variables and signed constraint,
objective, and right-hand-side data.  Positive linear programming
conventionally refers to the packing/covering pair with nonnegative data
\cite{LubyNisan1993PositiveLinearProgramming}.  We use the packing form,
which maximizes a nonnegative objective subject to nonnegative
upper-bound constraints, and give an exact sparse conversion from a
general LP threshold decision to a positive packing threshold.

Let $A\in\mathbb{Q}^{m\times n}$, $b\in\mathbb{Q}^m$,
$c\in\mathbb{Q}^n$, and $\tau\in\mathbb{Q}$.  The source question is
\begin{equation}
  \label{eq:lp-threshold}
  \exists x\in\R^n
  \quad
  Ax\leq b,
  \qquad
  c^\top x\geq\tau.
\end{equation}
Appending the inequality $-c^\top x\leq-\tau$ turns this question into
rational linear feasibility.  A polynomial-bit box containing a
feasible point can be computed directly from the rational input.  The
algebraic core of the reduction therefore considers arbitrary real
$A\in\R^{m\times n}$ and $b\in\R^m$ together with $U>0$, and handles
\begin{equation}
  \label{eq:boxed-system-intro}
  Ax\leq b,\qquad -U\one\leq x\leq U\one.
\end{equation}
The entries of $A$ and $b$ may have either sign.  The target is a
positive packing linear program
\begin{equation}
  \label{eq:packing-form-intro}
  \max\{e^\top z:z\geq0,\ Pz\leq d\},
\end{equation}
in which $P$, $d$, and $e$ are entrywise nonnegative.  Our goal is not
to solve either program, but to convert
\eqref{eq:boxed-system-intro} to an exact objective-threshold question
for \eqref{eq:packing-form-intro} while preserving sparsity.  Combining
this core with the one-row objective reduction and the polynomial-bit
box gives the claimed reduction for the general rational question
\eqref{eq:lp-threshold}.

For rational input, general linear programming is polynomial-time
solvable by the ellipsoid and interior-point frameworks
\cite{Khachiyan1979PolynomialLP,Karmarkar1984NewPolynomialTimeLP}.
That bit-complexity statement is distinct from the unit-cost
construction bound proved here.  Exact positive linear programming
also has a parallel-complexity barrier: Trevisan and Xhafa proved that
exact positive-LP solution is P-complete, so an NC algorithm would
imply $\mathrm{P}=\mathrm{NC}$
\cite{TrevisanXhafa1998ParallelComplexity}.

Much of the algorithmic literature therefore concerns approximation.
For a packing program with optimum $\OPT$, a
$(1-\varepsilon)$-approximate solution is a feasible point of value at
least $(1-\varepsilon)\OPT$.  Luby and Nisan initiated
width-independent parallel approximation for positive linear programs
\cite{LubyNisan1993PositiveLinearProgramming}; Plotkin, Shmoys, and
Tardos developed the Lagrangian-relaxation framework for fractional
packing and covering \cite{PlotkinShmoysTardos1995FractionalPackingCovering}.
For an explicit constraint matrix with $N$ nonzeros, $r$ rows, and $s$
columns, Koufogiannakis and Young gave, with high probability, feasible
primal and dual solutions within a factor $1+\varepsilon$ of optimum in
\[
  O\!\left(N+\frac{(r+s)\log N}{\varepsilon^2}\right)
\]
time \cite{KoufogiannakisYoung2014ExplicitFractionalPackingCovering}.
Allen-Zhu and Orecchia subsequently obtained
$\widetilde O(N/\varepsilon)$ time for packing and
$\widetilde O(N/\varepsilon^{3/2})$ time for covering
\cite{AllenZhuOrecchia2019NearlyLinearPackingCovering}.

The direct motivation for asking about this conversion comes from
Kyng, Wang, and Zhang
\cite{KyngWangZhang2020PackingLPsHard}.  They give quantitative
reductions from conditioned least-squares problems to packing programs,
showing that sufficiently fast and accurate packing solvers would yield
faster algorithms for linear systems.  More precisely, for a
full-column-rank matrix $A\in\R^{m\times n}$ and a known bound
$K_A\geq\kappa(A)$, their sparse reduction turns a
$\widetilde O(N/\varepsilon^\alpha)$-time packing solver into a
constant-relative-residual least-squares solver running in
$\widetilde O(\nnz(A)(m^{3/2}K_A^3)^\alpha)$ time.  The condition-number
bound determines the solution scale and packing accuracy used for
residual decoding and iterative refinement.

At the center of our reduction is a simpler exact conversion from a
signed system inside a known box to positive linear programming.  For
rational input, a determinant bound supplies a suitable box of
polynomial encoding length, so no boundedness promise is imposed on the
source feasibility problem.  At a structural level, our construction
and that of Kyng, Wang, and Zhang both split signed variables and use an
objective to force bounding constraints to be tight.  The guarantee
here is otherwise different: it does not provide the quantitative
approximation-to-residual implication needed for the linear-system
application.  We make no novelty claim for the shared algebraic devices.
The purpose of this note is to give a self-contained exact formulation,
exact sparse accounting, and an explicit statement of what
near-threshold points do and do not imply.

We first state the boxed-feasibility primitive.

\begin{theorem}[Boxed linear feasibility to positive linear programming]
\label{thm:main}
Let $m,n$ be positive integers, let $A\in\R^{m\times n}$ be given in
sparse row form, let $b\in\R^m$, and let $U>0$.  In the exact sparse
arithmetic model, a deterministic algorithm constructs, using
$O(\nnz(A)+m+n)$ arithmetic operations and output writes, a packing
linear program
\[
  \max\{e^\top z:z\in\R_{\geq0}^{q},\ Pz\leq d\}
\]
and a threshold $\theta>0$ such that $P$, $d$, and $e$ are entrywise
nonnegative,
\[
  q\leq2n,\qquad \nnz(P)\leq\nnz(A)+2n,
\]
and
\[
  \exists x\in[-U,U]^n\text{ with }Ax\leq b
  \quad\Longleftrightarrow\quad
  \max\{e^\top z:z\geq0,\ Pz\leq d\}\geq\theta.
\]
In the nontrivial branch the output has $q=2n$, exactly $m+n$
constraints, $\theta=nU$, and
$\nnz(P)=\nnz(A)+2n$.  A boxed feasible point can be encoded as a
packing point of value $\theta$ in $O(n)$ operations; conversely, any
explicit feasible packing point of value at least $\theta$ decodes to
a boxed feasible point in $O(n)$ operations.
\end{theorem}

\begin{lemma}[A polynomial-bit box for rational feasibility]
\label{lem:rational-box}
Let $p,n$ be positive integers, let $B\in\mathbb{Q}^{p\times n}$, and
let $d\in\mathbb{Q}^p$.  Write the entries in lowest terms with positive
denominators.  For each row $i$, let $q_i$ be the product of the
denominators appearing in that row, including the denominator of $d_i$,
and multiply the row by $q_i$.  This gives an equivalent integral system
$\overline Bx\leq\overline d$.  Define
\[
  M=\max\bigl(\{1\}\cup
  \{|\overline B_{ij}|:i\in[p],j\in[n]\}\cup
  \{|\overline d_i|:i\in[p]\}\bigr),
  \qquad
  r=\min\{p,n\},
\]
and $U=(rM)^r$.  Then
\[
  \{x\in\R^n:Bx\leq d\}\neq\varnothing
  \quad\Longleftrightarrow\quad
  \{x\in[-U,U]^n:Bx\leq d\}\neq\varnothing.
\]
Moreover,
\[
  \log_2 U\leq r(\log_2 r+\log_2 M),
\]
so $U$ has polynomial encoding length in the binary encoding of
$(B,d)$.
\end{lemma}

\begin{proof}
Only the forward implication requires proof.  Write
$x=x^+-x^-$ and introduce slacks:
\[
  \overline Bx^+-\overline Bx^-+s=\overline d,
  \qquad x^+,x^-,s\geq0.
\]
If this standard-form system is feasible, it has a basic feasible
solution.  One way to see this is to minimize the sum of all variables:
the sublevel set below the value of any fixed feasible point is compact,
so an optimum exists, and an extreme optimum is basic.

Let $C$ be the integral basis matrix of such a solution, and suppose
that $k$ of its columns come from
$(\overline B,-\overline B)$ rather than from the slack identity
matrix.  A nonsingular basis cannot contain both
$\overline B_{*,j}$ and $-\overline B_{*,j}$, so
$k\leq\min\{p,n\}=r$.  Expanding along the $p-k$ slack columns reduces
$\det C$ to a $k\times k$ integral determinant.  Hence
$|\det C|\geq1$.  For any basic coordinate belonging to $x^+$ or
$x^-$, Cramer's rule and the same expansion bound the numerator by
\[
  k!M^k\leq(rM)^r=U.
\]
At most one of $x_j^+$ and $x_j^-$ is basic; every nonbasic variable is
zero.  Therefore $|x_j|\leq U$ for every $j$.  The binary length of
$q_i$ is at most the sum of the binary lengths of the denominators in
row $i$.  Hence the normalized integer data, $M$, and $U$ all have
encoding length polynomial in that of $(B,d)$.
\end{proof}

\begin{corollary}[Rational LP threshold decision to a positive packing threshold]
\label{cor:general-lp}
Let $A\in\mathbb{Q}^{m\times n}$, $b\in\mathbb{Q}^m$,
$c\in\mathbb{Q}^n$, and $\tau\in\mathbb{Q}$ be binary encoded, with
$m,n$ positive.  A deterministic polynomial-time algorithm constructs
a positive packing linear program and a threshold $\theta>0$ such that
\[
  \exists x\in\R^n
  \quad
  Ax\leq b,
  \qquad
  c^\top x\geq\tau
  \quad\Longleftrightarrow\quad
  \OPT_{\mathrm{pack}}\geq\theta.
\]
The output has at most $2n$ variables and at most $m+n+1$
constraints, its constraint matrix $P$ satisfies
\[
  \nnz(P)\leq\nnz(A)+\nnz(c)+2n,
\]
and every output coefficient has encoding length polynomial in the
input encoding length.  An explicit packing point of value at least
$\theta$ decodes to a source feasible point in $O(n)$ exact arithmetic
operations.
\end{corollary}

\begin{proof}
Append the objective row and write the resulting feasibility system as
\[
  Bx\leq d,
  \qquad
  B=
  \begin{pmatrix}
    A\\ -c^\top
  \end{pmatrix},
  \qquad
  d=
  \begin{pmatrix}
    b\\ -\tau
  \end{pmatrix}.
\]
Normalize each row using the denominator product specified in
Lemma~\ref{lem:rational-box}, and compute $U$ from
that lemma.  Applying \Cref{thm:main} to the normalized
system inside $[-U,U]^n$ proves the equivalence and the dimension,
sparsity, and witness claims.  Row normalization preserves the zero
pattern.  Its coefficient lengths are polynomial in the input length,
as is $\log U$; the shifted coefficients, shifted right-hand sides, box
bounds, and threshold therefore also have polynomial encoding length.
All normalization, integer arithmetic, and sparse-output construction
can be performed in polynomial bit time.
\end{proof}

\paragraph{Proof idea.}
For rational input, first write the objective threshold as one
additional row, clear denominators, and use
Lemma~\ref{lem:rational-box} to compute a polynomial-bit coordinate bound.
For the sparse boxed-feasibility core, write $x_j=y_j-s_j$ with
$y_j,s_j\geq0$ and intend to enforce
$y_j+s_j=U$.  For row $i$, put
\[
  R_i=\sum_j|A_{ij}|,
  \qquad p^+_{ij}=A_{ij}+|A_{ij}|,
  \qquad p^-_{ij}=-A_{ij}+|A_{ij}|.
\]
The shifted row
\[
  \sum_jp^+_{ij}y_j+\sum_jp^-_{ij}s_j
  \leq b_i+UR_i
\]
has nonnegative coefficients and, on the intended slice, is exactly
$A_i x\leq b_i$.  The packing constraints $y_j+s_j\leq U$ and the
objective $\sum_j(y_j+s_j)$ with threshold $nU$ force every one of
these inequalities to be tight.  If some shifted right-hand side is
negative, the corresponding original row is infeasible throughout the
box and the algorithm returns a fixed rejecting instance.

\paragraph{Scope.}
Removing the supplied box uses the finite binary encoding of rational
data.  For arbitrary exact-real data, \Cref{thm:main} still requires
$U$ as input.  The computed $U$ may have exponential magnitude; only
its encoding length is polynomial.  Thus Corollary~\ref{cor:general-lp} is a
polynomial bit-time reduction, not a strongly polynomial result.  The
$O(\nnz(A)+m+n)$ bound in \Cref{thm:main} counts exact unit-cost
arithmetic operations and output writes after $U$ is available, and it
does not include the cost of solving the output LP.  The construction
is an exact-threshold reduction, not a promise-gap reduction.  We prove in
\Cref{sec:approximation-boundary} that near-threshold packing points
decode to additively near-feasible boxed points, but also that
infeasible instances can have packing optimum arbitrarily close to
$nU$.  Thus a value-only threshold oracle does not supply a witness,
and a standard multiplicative approximation does not decide arbitrary
boxed feasibility.

The remainder of the paper is organized as follows.
\Cref{sec:setup} records the algebraic identity behind the shift.
\Cref{sec:reduction} gives the algorithm and proves the feasibility
primitive.  \Cref{sec:approximation-boundary} gives the
approximate decoding bound and the no-gap example.

\section{Setup}
\label{sec:setup}

All vector inequalities are coordinatewise.  For a positive integer
$k$, write $[k]=\{1,\ldots,k\}$, and let $\one$ denote an all-ones
vector of the dimension determined by context.  Empty sums are zero.

\begin{definition}[Positive packing threshold instance]
\label{def:packing-threshold}
A packing linear program, the maximization form of a positive linear
program, has the form
\[
  \max\{e^\top z:z\in\R_{\geq0}^{q},\ Pz\leq d\},
\]
where $P$, $d$, and $e$ are entrywise nonnegative.  A packing
threshold instance also specifies $\theta>0$ and accepts if the
maximum is at least $\theta$.
\end{definition}

\paragraph{Sparse arithmetic model.}
The input matrix is available in sparse row form.  Reading and writing
a stored coefficient, adding or comparing real numbers, taking an
absolute value, and multiplying by $U$ each cost one operation.  These
are exact operations.  The cost of writing the sparse output and its
dense right-hand side and objective vectors is included.

For $i\in[m]$ and $j\in[n]$, define
\begin{equation}
  \label{eq:shift-data}
  R_i=\sum_{j=1}^n|A_{ij}|,
  \qquad
  p^+_{ij}=A_{ij}+|A_{ij}|,
  \qquad
  p^-_{ij}=-A_{ij}+|A_{ij}|,
  \qquad
  c_i=b_i+UR_i.
\end{equation}
The quantities $R_i,p^+_{ij},p^-_{ij}$ are nonnegative.  The
right-hand side $c_i$ need not be.

\begin{observation}[Shift identity and tight slice]
\label{obs:shift-tight}
For all $y,s\in\R_{\geq0}^n$ and $i\in[m]$,
\begin{equation}
  \label{eq:shift-identity}
  \sum_{j=1}^n p^+_{ij}y_j+
  \sum_{j=1}^n p^-_{ij}s_j
  =
  \sum_{j=1}^n A_{ij}(y_j-s_j)+
  \sum_{j=1}^n |A_{ij}|(y_j+s_j).
\end{equation}
If $y_j+s_j\leq U$ for every $j$ and
$\sum_j(y_j+s_j)\geq nU$, then $y_j+s_j=U$ for every
$j$.
\end{observation}

\begin{proof}
Expanding the definitions of $p^+_{ij}$ and $p^-_{ij}$ gives
\eqref{eq:shift-identity}.  For the second assertion, set
$\delta_j=U-y_j-s_j\geq0$.  The objective inequality says
$\sum_j\delta_j\leq0$, so every nonnegative $\delta_j$ is zero.
\end{proof}

The tight slice is in bijection with the box.  If $x\in[-U,U]^n$,
then
\begin{equation}
  \label{eq:encode}
  y_j=\frac{U+x_j}{2},
  \qquad
  s_j=\frac{U-x_j}{2}
\end{equation}
is nonnegative, $y_j+s_j=U$, and $y_j-s_j=x_j$.
Conversely, nonnegative $y_j,s_j$ with sum $U$ give
$x_j=y_j-s_j\in[-U,U]$.

\section{The reduction to positive linear programming}
\label{sec:reduction}

We first describe the algorithm.  Coefficients equal to zero are not
stored.

\paragraph{Algorithm \textsc{BoxToPacking}.}
Given $A,b,U$, perform the following steps.
\begin{enumerate}[label=\arabic*.,leftmargin=2em]
  \item Compute $R_i$ and $c_i$ from \eqref{eq:shift-data} for every
  $i\in[m]$.

  \item If $c_i<0$ for some $i$, return the one-variable packing
  program
  \[
    \max\{0\cdot z:z\geq0,\ z\leq0\}
  \]
  with threshold $\theta=1$.

  \item Otherwise, use variables
  \[
    z=(y_1,\ldots,y_n,s_1,\ldots,s_n)\in\R_{\geq0}^{2n},
  \]
  set the objective vector to the all-ones vector, and set
  $\theta=nU$.

  \item For each $i\in[m]$, create the shifted row
  \begin{equation}
    \label{eq:shifted-row}
    \sum_{j=1}^n p^+_{ij}y_j+
    \sum_{j=1}^n p^-_{ij}s_j\leq c_i.
  \end{equation}

  \item For every $j\in[n]$, create the box row
  \begin{equation}
    \label{eq:box-row}
    y_j+s_j\leq U.
  \end{equation}
  Return these $m+n$ constraints, the objective, and the threshold.
\end{enumerate}

The exceptional branch is justified by the following elementary
certificate.

\begin{lemma}[A negative shifted right-hand side certifies infeasibility]
\label{lem:negative-shift}
If $c_i<0$ for some $i\in[m]$, then the boxed system is infeasible.
\end{lemma}

\begin{proof}
For every $x\in[-U,U]^n$,
\[
  (Ax)_i
  \geq-\sum_{j=1}^n|A_{ij}|\,|x_j|
  \geq-UR_i.
\]
The inequality $c_i=b_i+UR_i<0$ is equivalent to
$b_i<-UR_i$.  Hence $(Ax)_i>b_i$ throughout the box.
\end{proof}

We next prove the exact equivalence in the main branch.

\begin{lemma}[Threshold equivalence and witness conversion]
\label{lem:threshold-equivalence}
Assume $c_i\geq0$ for every $i\in[m]$.  The boxed system is feasible
if and only if the packing program defined by
\eqref{eq:shifted-row} and \eqref{eq:box-row} has optimum at least
$nU$.  The maps $x\mapsto(y,s)$ from \eqref{eq:encode} and
$(y,s)\mapsto y-s$ convert explicit witnesses in $O(n)$ operations.
\end{lemma}

\begin{proof}
Suppose $x\in[-U,U]^n$ and $Ax\leq b$.  Define $y,s$ by
\eqref{eq:encode}.  They satisfy every box row with equality.  By
\eqref{eq:shift-identity}, for each $i\in[m]$,
\[
  \sum_jp^+_{ij}y_j+\sum_jp^-_{ij}s_j
  =(Ax)_i+UR_i
  \leq b_i+UR_i=c_i.
\]
Thus $(y,s)$ is packing-feasible and has objective $nU$.

Conversely, the origin is packing-feasible because every $c_i$ is
nonnegative.  The box rows make the feasible region compact, so the
maximum is attained.  Let $(y,s)$ be feasible with objective at least
$nU$.  The box rows also bound that objective by $nU$; hence
Observation~\ref{obs:shift-tight} gives $y_j+s_j=U$ for every $j$.  Put
$x=y-s$.  Then $x\in[-U,U]^n$, and
\eqref{eq:shift-identity} gives
\[
  (Ax)_i+UR_i
  =\sum_jp^+_{ij}y_j+\sum_jp^-_{ij}s_j
  \leq c_i=b_i+UR_i.
\]
Therefore $Ax\leq b$.  Both conversions require constant work per
coordinate.
\end{proof}

\begin{proof}[Proof of \Cref{thm:main}]
Run \textsc{BoxToPacking}.  If the exceptional branch is taken, its
packing optimum is $0<1=\theta$, while
Lemma~\ref{lem:negative-shift} says that the boxed system is infeasible.
The branch has one variable, one constraint, and one matrix nonzero;
because $n\geq1$,
\[
  1\leq\nnz(A)+2n.
\]

In the main branch, every coefficient in
\eqref{eq:shifted-row} is nonnegative, every $c_i$ is nonnegative, and
the box rows, objective, and threshold are positive or nonnegative as
required.  Thus the output is a packing threshold instance.
Lemma~\ref{lem:threshold-equivalence} proves the exact equivalence and the
witness statements.

It remains to count the output and the construction work.  If
$A_{ij}>0$, then $p^+_{ij}=2A_{ij}$ and $p^-_{ij}=0$; if
$A_{ij}<0$, then $p^+_{ij}=0$ and $p^-_{ij}=-2A_{ij}$.  Consequently,
each stored nonzero of $A$ produces exactly one nonzero in the shifted
rows.  The $n$ box rows contribute two nonzeros each, so
\[
  \nnz(P)=\nnz(A)+2n
\]
in the main branch.  A constant number of sparse scans computes the
row radii and writes the shifted coefficients.  Initializing the $m$
right-hand sides, the $2n$ objective entries, and the $n$ box rows
costs $O(m+n)$ more operations and writes.  This proves the claimed
$O(\nnz(A)+m+n)$ bound.
\end{proof}

\section{The approximation boundary}
\label{sec:approximation-boundary}

Although the exact reduction has no promise gap, its feasible points
have a simple quantitative interpretation below the threshold.

\begin{proposition}[Decoding a near-threshold point]
\label{prop:approximate-decoding}
Consider a main-branch instance and any packing-feasible point
$(y,s)$.  Define
\[
  x=y-s,
  \qquad
  \delta_j=U-y_j-s_j,
  \qquad
  \Delta=\sum_{j=1}^n\delta_j
  =nU-\sum_{j=1}^n(y_j+s_j).
\]
Then $x\in[-U,U]^n$, and for every $i\in[m]$,
\begin{equation}
  \label{eq:approximate-row-bound}
  (Ax)_i
  \leq b_i+\sum_{j=1}^n|A_{ij}|\delta_j
  \leq b_i+\lVert A_{i,*}\rVert_\infty\Delta.
\end{equation}
In particular, objective value at least $(1-\varepsilon)nU$ implies
\[
  (Ax)_i\leq b_i+
  \varepsilon nU\lVert A_{i,*}\rVert_\infty
  \qquad(i\in[m]).
\]
\end{proposition}

\begin{proof}
Nonnegativity and the box rows give
$|x_j|=|y_j-s_j|\leq y_j+s_j\leq U$.  Rearranging the shifted-row
inequality with \eqref{eq:shift-identity} yields
\[
  (Ax)_i
  \leq b_i+UR_i-\sum_j|A_{ij}|(y_j+s_j)
  =b_i+\sum_j|A_{ij}|\delta_j.
\]
The second inequality in \eqref{eq:approximate-row-bound} follows
from $\delta_j\geq0$ and $\sum_j\delta_j=\Delta$.  The final statement
uses $\Delta\leq\varepsilon nU$.
\end{proof}

The additive bound in Proposition~\ref{prop:approximate-decoding} does not yield
an exact feasibility decision without a separate promise.  The next
example shows that this is unavoidable for the present construction.

\begin{proposition}[No uniform objective gap]
\label{prop:no-gap}
For every $\delta\in(0,1)$, there is an infeasible boxed instance with
$m=2$, $n=1$, and $U=1$ whose output takes the main branch and has
\[
  \OPT=\theta-\frac{\delta}{2}.
\]
Consequently, the separation between infeasible and feasible outputs
can be arbitrarily small, both additively and multiplicatively.
\end{proposition}

\begin{proof}
Take
\[
  A=\begin{pmatrix}1\\-1\end{pmatrix},
  \qquad
  b=\begin{pmatrix}0\\-\delta\end{pmatrix},
  \qquad U=1.
\]
The boxed system requires $x\leq0$ and $x\geq\delta$, so it is
infeasible.  Both row radii are $1$, and the shifted right-hand sides
are $1$ and $1-\delta$, so the main branch is taken.  Its variables
are $y,s\geq0$, its threshold is $\theta=1$, and its constraints are
\[
  2y\leq1,
  \qquad
  2s\leq1-\delta,
  \qquad
  y+s\leq1.
\]
The first two bounds are simultaneously attainable and imply
\[
  \OPT=\frac12+\frac{1-\delta}{2}
  =1-\frac{\delta}{2}.
\]
Letting $\delta$ tend to zero proves the claim.
\end{proof}

Thus the exact threshold equivalence alone cannot support an
approximation-preserving reduction.  Any such extension would need an
additional promise or a different construction that creates and tracks
a quantitative gap.

\bibliographystyle{alpha}
\bibliography{references}

\end{document}
